\specialsection{Обзор литературы}

В данной работе рассматривается задача адаптивного выбора избыточности для уже
записанных данных в распределённом хранилище. Поэтому в обзоре важны не только
сами методы отказоустойчивого хранения, но и то, как существующие системы
учитывают активность данных и выполняют переходы между схемами хранения.

\subsection*{Репликация и коды стираний}

Классический выбор в распределённых системах хранения лежит между репликацией
и кодами стираний. Репликация проста: если одна реплика недоступна, данные можно
прочитать из другой реплики, а восстановление потерянного блока сводится к
копированию. Однако трёхкратная репликация требует хранить втрое больше данных,
что плохо масштабируется для крупных кластеров.

Коды стираний решают противоположную задачу: они уменьшают объём избыточности,
но повышают стоимость чтения и восстановления при отказах. В обзоре современных
подходов этот компромисс сформулирован так:
\enquote{erasure coding incurs high performance overhead in repair and updates}
\cite{shen2024erasurecoding}. Классическим семейством таких кодов являются коды
Рида--Соломона \cite{reed1960polynomial}. Практический опыт Azure и Facebook f4
показывает, что переход от репликации к кодам стираний даёт существенную
экономию места, но требует отдельно решать проблему дорогого восстановления и
чтения при недоступности блоков \cite{huang2012azure,muralidhar2014f4}.
Кроме числа читаемых блоков, важна и вычислительная цена кодирования:
сравнение открытых библиотек кодов стираний для систем хранения показывает, что
выбор реализации заметно влияет на пропускную способность кодирования и
декодирования. В работе о производительности кодирования отдельно отмечено,
что \enquote{parameter selection, especially as it concerns memory, has a
significant impact} \cite{plank2009performance}.

Ограничение этих подходов состоит в том, что они хорошо описывают крайние
режимы: либо данные хранятся как набор копий, либо как кодовая схема. Для
данных с меняющейся активностью этого недостаточно: горячим данным выгодны
дешёвые чтения, а холодным --- компактное хранение. Следовательно, возникает
потребность в промежуточных состояниях и переходах между ними.

\subsection*{Локальное восстановление}

Одна из попыток смягчить недостатки кодов стираний связана с локально
восстанавливаемыми кодами (locally recoverable codes, LRC). Их идея состоит
в том, чтобы добавить локальные проверочные блоки и тем самым уменьшить число
блоков, которые нужно читать при одиночном отказе. Теоретическая постановка
локальности рассмотрена в \cite{gopalan2012lrc}, а практическое использование
локально восстанавливаемых кодов в Azure описано в \cite{huang2012azure}.
Практический смысл этой идеи авторы Azure описывают через уменьшение числа
читаемых блоков: \enquote{LRC reduces the number of erasure coding fragments
that need to be read} \cite{huang2012azure}. Схожая мотивация есть в Xorbas:
авторы показывают, что коды Рида--Соломона
экономят место, но создают высокую стоимость восстановления, и предлагают коды
с локальностью для снижения числа операций чтения и сетевого трафика
\cite{sathiamoorthy2013xoring}. Для практической системы авторы Xorbas
сообщают о \enquote{reduction of approximately 2x on the repair disk I/O}
\cite{sathiamoorthy2013xoring}.

Локальность, однако, не устраняет саму проблему выбора схемы. Она добавляет
новую точку в пространстве компромиссов: восстановление становится дешевле, но
появляются дополнительные параметры кода и требования к размещению блоков.
Работа о широких LRC подчёркивает, что при увеличении ширины кодовой полосы
возникают дополнительные вопросы надёжности и размещения: \enquote{wide LRC
reliability is a subtle phenomenon} \cite{kadekodi2023widelrc}.
Для данной работы это важно как ограничение: локальные проверочные блоки
полезны не только как часть фиксированного кода, но и как возможное
промежуточное состояние между репликацией и более компактными схемами.

\subsection*{Температура данных и уровни хранения}

Многие практические системы исходят из того, что данные имеют разную
активность. Небольшая часть объектов читается часто, а значительная часть
остаётся холодной. В работе \cite{levandoski2013hotcold} рассматривается
идентификация горячих и холодных данных и используются методы сглаживания
оценок частоты обращений. Для систем хранения это означает, что форма
избыточности должна зависеть не только от требуемой надёжности, но и от
ожидаемой нагрузки чтения.

Ряд систем использует это наблюдение на уровне классов или уровней хранения.
Facebook f4 ориентирован на тёплые неизменяемые крупные двоичные объекты, для
которых репликация становится слишком дорогой по месту; сами авторы называют
f4 \enquote{a warm BLOB storage system} \cite{muralidhar2014f4}.
ER-Store классифицирует части таблиц по температуре и выбирает между
репликацией и кодами стираний для разных классов данных \cite{li2021erstore}.
ELECT использует температуру в хранилище ключ--значение на основе дерева
лог-структурного слияния (LSM-дерева) и избирательно переводит данные между
репликацией и кодированием \cite{ren2024elect}. Близкий подход реализует
CBase-EC, который для распределённой базы данных с разделением чтения и записи
определяет горячие и холодные таблицы, а затем выполняет двунаправленные
переходы между трёхкратной репликацией и LRC-схемами с сохранением целевой
производительности транзакций. В этой работе явно подчёркивается, что нужны
не только переходы охлаждения, но и обратные переходы: \enquote{not only
necessary to design hot-to-cold approaches to reduce storage overhead, but also
cold-to-hot methods} \cite{yao2021cbaseec}. Метод HSM расширяет это
направление, учитывая не только температуру данных, но и текущую утилизацию
дискового пространства кластера: при свободном месте предпочтение отдаётся
схемам с дешёвым чтением, а при дефиците --- более компактным схемам
\cite{song2024hsm}.

Ограничение таких решений состоит в том, что адаптация обычно задаётся как
выбор уровня или класса хранения: горячие данные остаются ближе к репликации,
холодные переводятся ближе к кодированию. Такой подход хорошо работает для
систем с естественными уровнями, но он не отвечает на вопрос, какой именно
локальный переход выгоден для конкретной уже записанной схемы. Кроме того,
обратные переходы при повторном нагреве данных обычно не рассматриваются как
часть единого пространства состояний.

\subsection*{Гибридные схемы}

Между чистой репликацией и чистым кодом стираний существуют гибридные подходы.
HyRES объединяет репликацию, кодирование со свойством максимального
расстояния по столбцам и элементы локального восстановления, а затем
сравнивает такие схемы по стоимости хранения, вероятности потери файла и
трафику восстановления \cite{lucani2025hyres}. Cocytus, реализованный для
in-memory хранилища ключ--значение, идёт в ту же сторону на уровне системы:
его цель описывается как \enquote{reducing the memory consumption for replicas
while keeping similar performance} \cite{zhang2016cocytus}; небольшие и часто
обновляемые объекты он реплицирует, а крупные значения
кодирует, добиваясь снижения расхода памяти на 33--46\,\% по сравнению с
полным репликационным резервированием при сопоставимой задержке
\cite{zhang2016cocytus}. Эта линия работ важна тем, что показывает: реплики
и кодовые блоки можно рассматривать не как взаимоисключающие решения, а как
компоненты одной схемы избыточности.

Однако HyRES в первую очередь анализирует свойства семейства гибридных схем, а
не механизм их динамического изменения в работающем кластере. Cocytus, ER-Store
и ELECT, в свою очередь, используют гибридность на уровне политики хранения,
но не строят последовательность малых обратимых переходов внутри схемы одного
объекта \cite{zhang2016cocytus,li2021erstore,ren2024elect}. Из-за этого
адаптация получается
достаточно грубой: перенос данных между уровнями или классами хранения может
резко менять расход места и нагрузку чтения, а также создавать локальный
дисбаланс использования ресурсов. Поэтому остаётся открытым архитектурный
вопрос: как представить реплики, локальные и глобальные проверочные блоки
в одной модели, чтобы система могла постепенно менять их соотношение, а не
переключать данные между крупными режимами хранения.

\subsection*{Изменение схемы уже записанных данных}

Если схема хранения меняется после записи данных, возникает стоимость
перекодирования. В общем случае нужно прочитать существующие блоки,
сформировать новую избыточность и удалить старую. Конвертируемые коды
формализуют задачу изменения параметров кодирования и минимизации числа
прочитанных блоков при таком переходе. Авторы вводят это как
\enquote{code conversion---the process of converting data encoded with an
initial code} \cite{maturana2022convertible}.
Для локально восстанавливаемых кодов похожая задача рассматривается в режиме
объединения: \enquote{code conversions for locally repairable codes in the
merge regime} \cite{kong2024lrcconvertible}.

Системные работы по адаптивной избыточности предлагают разные ответы на
вопрос, когда и как менять схему. HeART наблюдает за фактической надёжностью
дисков различных групп и подбирает схему избыточности под текущую интенсивность
отказов: для более надёжных групп допустимы более компактные коды
\cite{kadekodi2019heart}. Мотивация HeART состоит в отказе от режима
\enquote{one-scheme-for-all fashion} \cite{kadekodi2019heart}. Продолжение этой линии, PACEMAKER, обращает
внимание на то, что массовые переходы между схемами могут перегружать
полосу пропускания кластера, и предлагает оркестратор переходов, который
заранее организует размещение блоков и сглаживает фоновую нагрузку на
ввод-вывод. Масштаб этой проблемы авторы формулируют жёстко:
\enquote{transition IO under existing approaches can consume 100\% cluster IO
continuously for several weeks} \cite{kadekodi2020pacemaker}. Tiger развивает оба направления,
снимая ограничения на размещение, которые требовались в предшествующих
системах с адаптивной избыточностью; авторы формулируют это как
\enquote{avoids adoption-blocking placement constraints} \cite{kadekodi2022tiger}. Отдельно
работа EAR показывает, что уже на этапе репликации стоит учитывать будущее
перекодирование: размещение реплик с учётом будущей сборки в коды стираний
позволяет избежать межстоечных передач и переноса данных после
кодирования. Причина в том, что \enquote{random replication does not take into
account erasure coding} \cite{li2017ear}.

Наиболее близкой системной работой является Morph \cite{kim2024morph}. Morph
исходит из наблюдения, что в реальных кластерных файловых системах данные
проходят жизненные фазы: сначала они часто хранятся с трёхкратной репликацией,
затем переводятся в более компактные схемы; авторы описывают это как необходимость
\enquote{change (transcode) the redundancy schemes of those files} \cite{kim2024morph}.
Авторы показывают, что операции
перекодирования могут создавать значительную нагрузку на кластер, и предлагают
схемы и политики размещения, уменьшающие стоимость таких переходов.

Ограничение Morph, HeART и PACEMAKER относительно постановки данной работы
состоит в том, что жизненный цикл и целевые переходы в значительной степени
задаются заранее: новые данные охлаждаются и переводятся в более компактные
схемы либо схема выбирается по наблюдаемой надёжности устройств. В данной
работе целевое состояние не фиксируется заранее. Планировщик должен выбирать
переход из текущего состояния на основе текущей активности данных, занятого
места и ожидаемой стоимости чтения. Это делает задачу не только задачей
эффективного перекодирования, но и задачей выбора полезного перехода.

\subsection*{Выводы по обзору литературы}

Рассмотренные работы показывают, что отдельные части задачи хорошо изучены:
репликация обеспечивает дешёвые чтения, коды стираний экономят место,
локально восстанавливаемые коды уменьшают стоимость одиночных отказов,
температурные политики позволяют разделять горячие и холодные данные, а
конвертируемые коды снижают стоимость заранее заданных переходов. При этом
каждое направление оставляет ограничение, важное для данной работы.

Во-первых, статический выбор между репликацией и кодами стираний плохо
подходит для данных, активность которых меняется. Во-вторых, температурные
системы часто выбирают класс или уровень хранения, но не строят общее
пространство обратимых переходов для одной схемы. В-третьих, работы о
конвертируемых кодах и Morph уменьшают стоимость перехода, но не решают в
общем виде вопрос, какой переход следует выполнить сейчас. Поэтому в данной
работе предлагается рассматривать схему хранения как изменяемое состояние,
состоящее из реплик, локальных и глобальных проверочных блоков, а выбор
переходов поручать планировщику, учитывающему место, активность данных и
ожидаемую стоимость чтения.

\pagebreak
